\documentclass[a4paper,10pt]{article}
\usepackage[utf8]{inputenc}
\usepackage[spanish,es-tabla]{babel}

\usepackage[a4paper,headheight=43pt,
            scale={0.8,0.8},hoffset=0.5cm,voffset=1cm]{geometry}

\usepackage{verbatim}
\usepackage{comment}

% paquetes matemáticos básicos
\usepackage{amsmath}
\usepackage{amsfonts}
\usepackage{amssymb}
\usepackage{xfrac}
\usepackage{siunitx}
\usepackage{enumitem,amssymb}
\newlist{todolist}{itemize}{2}
\setlist[todolist]{label=$\square$}

% tikz
\usepackage{tikz}

% Código
\usepackage{listings}

\definecolor{codegreen}{rgb}{0,0.6,0}
\definecolor{codegray}{rgb}{0.5,0.5,0.5}
\definecolor{codepurple}{rgb}{0.58,0,0.82}
\definecolor{backcolour}{rgb}{0.95,0.95,0.92}

\lstdefinestyle{mystyle}{
    backgroundcolor=\color{backcolour},   
    commentstyle=\color{codegreen},
    keywordstyle=\color{magenta},
    numberstyle=\tiny\color{codegray},
    stringstyle=\color{codepurple},
    basicstyle=\ttfamily\footnotesize,
    breakatwhitespace=false,         
    breaklines=true,                 
    captionpos=b,                    
    keepspaces=true,                 
    numbers=left,                    
    numbersep=5pt,                  
    showspaces=false,                
    showstringspaces=false,
    showtabs=false,                  
    tabsize=2
}

\lstset{style=mystyle}

% gráficos
\usepackage{graphicx}
\usepackage{float}
\usepackage{caption}
\usepackage{subfig}
\usepackage{listings}
\captionsetup[table]{font={small,sc},labelfont={small,sc}}

\usepackage[subrefformat=parens]{subcaption} % para colocar subfiguras con sus
                                             % respectivas referencias,
                                             % subrefformat es para que ponga
                                             % parentesis en las subref

\usepackage{appendix}
\usepackage[output-decimal-marker={,},per-mode=symbol-or-fraction,retain-unity-mantissa=true,detect-all,range-phrase={ -- },list-final-separator={ y }]{siunitx}
%output-decimal-marker define el separador decimal que se imprime
%per-mode define la forma que se utiliza para mostrar las unidades del denominador, en este caso elije segun el modo matematico que usa
%retain-unity-mantissa define si muestra o no la mantisa cuando es 1
%detect-all detecta la fuente que corresponde al entorno donde se colocan las unidades
%range-phrase define la palabra para los rangos de numeros
%list-final-separator define la palabra final de las listas

% paquete para insertar hipervínculos (sólo para el encabezado)
\usepackage[colorlinks=true,urlcolor=blue,
            linkcolor=black,linktoc=none,citecolor=black]{hyperref}

\usepackage[twoside]{fancyhdr}		% formato de página
\usepackage{lastpage}
% Configuración del encabezado y pie de página
\fancyhf{} % Limpia todos los campos de encabezado y pie de página
%\fancyhead[L]{\textbf{TB070 Dispositivos Semiconductores}}
\fancyfoot{} % Limpia el pie de página para personalizarlo a continuación

\fancyfoot[LE,RO]{\thepage/\pageref{LastPage}}
\fancyfoot[LE,RO]{\thepage/\pageref{LastPage}}

%*** redefino el comportamiento de los apendices
\renewcommand{\appendixname}{Apéndices}
\renewcommand{\appendixtocname}{Apéndices}
\renewcommand{\appendixpagename}{Apéndices}
%\renewcommand{\tablename}{Tabla}
%\renewcommand\spanishtablename{Tabla}
%\renewcommand{\figurename}{Figura}
%***
\pagestyle{fancy}

\usepackage[dvipsnames]{xcolor}
%\usepackage{todonotes}		% ToDo List
\usepackage{hyperref}		% formato de página
\usepackage{tcolorbox}

\DeclareSIUnit\angstrom{\text {Å}}

% \newif\ifturnotarde
% %     \turnotardetrue
% \newif\ifturnonoche
%     \turnonochetrue

\begin{document}

% Defino el formato del encabezado
%\pagestyle{fancy}
% \lhead{\includegraphics[width=0.3\textwidth]{lhead}}
% \chead{
%   {\scshape Dispositivos Semiconductores}\\
%   {\footnotesize 1\textsuperscript{er} cuatrimestre de 2023}
% }
%\rhead{\includegraphics[width=0.1\textwidth]{rhead}}
\lfoot{}
\cfoot{\thepage}
\rfoot{}
\begin{titlepage}

\thispagestyle{empty}

\begin{center}
    \includegraphics[scale=0.3]{img/fiuba.pdf}\\
    \large{\textsc{Universidad de Buenos Aires}}\\
    \large{\textsc{Facultad de Ingeniería}}\\
    % Modificar año y cuatrimestre
    \small{Año 2026 - 2\textsuperscript{do} cuatrimestre}
\end{center}

\vfill
% Redefino la numeracion de las subsecciones
\renewcommand{\thesubsection}{\arabic{subsection}}

\begin{center}
	{\LARGE \underline{\textbf{Plan de Trabajo}}}\\[0.5cm]
	{\LARGE \textbf{Tesis de Grado de Ingeniería Electrónica}}\\[0.5cm] 
    {\LARGE \textbf{Incorporación de redundancia para sistemas móviles autónomos}}\\[0.5cm]
\end{center}

\begin{table}[H]
    \centering
    \begin{tabular}{|l|l|l|}
    \hline
    \multicolumn{3}{|c|}{\textbf{Alumno}}    \\ \hline
    Juan Ignacio Biancuzzo & 106005 & jbiancuzzo@fi.uba.ar \\ \hline\hline
    \multicolumn{3}{|c|}{\textbf{Director}}    \\ \hline
    Ignacio Mas &  & imas@fi.uba.ar \\ \hline
    \end{tabular}
\end{table}

\vfill

\end{titlepage}

\tableofcontents
\newpage

\section{Resumen}
\begin{comment}
    Breve descripción de como esta estructurado la propuesta 
\end{comment}

\section{Palabras claves}
\begin{comment}
    Todas las palabras claves que se utilicen 
        -> Ver de utilizar el package de abreviaciones que asi enunciamos de esa forma cada nombre importante
        -> Ir anotandolas mientras escribo
\end{comment}

\section{Introducción}
\begin{comment}
    Hacerlo al final, este debería describir todo la problematica del problema
\end{comment}


\section{Problemática}
\begin{comment}
    Describir porque se esta planteando este proyecto en especifico
\end{comment}


\section{Objetivos}
\begin{comment}
    Plantear en terminos generales que es lo que se busca realizar, y en particular como estará conformado
\end{comment}

\subsection{General}

\subsection{Especifico}


\section{Métodos}
\begin{comment}
    Describir como se realizará todo lo dicho
\end{comment}


\section{Experimentación y validación}
\begin{comment}
    Escribir que pruebas se necesitan realizar para poder validad el correcto funcionamiento de la tesis
\end{comment}


\section{Plan de trabajo}
\begin{comment}
    Hacer una tabla dividida por meses, y actividades, la idea sería expresar un cronograma de actividades
\end{comment}


\section{Referencias}
\renewcommand{\section}[2]{}
\bibliographystyle{IEEEtran}
\bibliography{referencias.bib}

\end{document}
